\begin{table*}[t]
\centering
\scriptsize
\renewcommand{\arraystretch}{0.9}
\begin{tabularx}{\linewidth}{>{\raggedright\arraybackslash}p{0.19\linewidth}>{\raggedright\arraybackslash}p{0.24\linewidth}>{\raggedright\arraybackslash}p{0.23\linewidth}X}
\toprule
\textbf{Apparent conclusion} & \textbf{What the framework changes} & \textbf{Non-trivial finding} & \textbf{Implication} \\
\midrule
``Method A is more robust to noise than method B.'' & Aligns noise type, rank position, gold-evidence availability, and intervention point. & Some methods solve generator distraction while others solve retriever defects; their scores are not competing estimates of the same property. & Report noise source and intervention point before ranking methods. \\
``Defense X lowers attack success.'' & Separates prevention, demotion, detection, traceback, and recovery. & A defense can improve traceback without lowering user-facing attack success, or certify bounded corruption without handling prompt-bearing documents. & Security papers should report outcome type, not only attack success rate. \\
``Calibration improves reliability.'' & Separates confidence in answer correctness from evidence sufficiency and action utility. & A calibrated answerer can still answer from insufficient evidence; an abstaining system can lower risk while reducing coverage. & Calibration metrics should be paired with evidence and action metrics. \\
``Benchmark Y is harder.'' & Compares evidence state, required action, cost budget, and metric target. & A benchmark can be hard because retrieval fails, because evidence conflicts, or because the correct action is refusal; these failures should not be pooled. & Leaderboards should expose diagnostic sub-scores instead of one robust-RAG score. \\
\bottomrule
\end{tabularx}
\caption{Distilled classes of non-trivial findings produced by the comparability framework. Main Table~\ref{tab:main-normalization-findings} gives the concrete paper-pair cases; this appendix table groups them into reusable interpretation rules.}
\label{tab:nontrivial-findings}
\end{table*}
